See \wikilink{s20--b001}{the character-group definition} and
\wikilink{s20--b006}{the finite evaluation pairing}.

\begin{proposition}\label{characters-abelianization}
For any group $G$, composition with abelianization gives a natural
group isomorphism
$\operatorname{Hom}(G^{\mathrm{ab}},\mathbb C^\times)
\cong\operatorname{Hom}(G,\mathbb C^\times)$.
\end{proposition}
\begin{proof}\label{characters-abelianization-proof}
The target is abelian, so every character factors uniquely through
\wikilink{s13--b006}{abelianization}. Composition respects pointwise
multiplication of characters, giving a group isomorphism.
\end{proof}
\begin{proposition}\label{finite-self-duality}
Every finite abelian group $A$ is isomorphic to its character group
$\widehat A=\operatorname{Hom}(A,\mathbb C^\times)$. The isomorphism
need not be canonical.
\end{proposition}
\begin{proof}\label{finite-self-duality-proof}
A character of a cyclic group of order $n$ is uniquely determined
by the image of a generator, an arbitrary $n$th root of unity.
These roots form a cyclic group of order $n$.
Characters of a finite direct sum are independently specified on
the summands, and their character group is the product of the
character groups of the summands. Apply finite abelian classification.
The resulting isomorphism uses choices of cyclic generators and
primitive roots of unity.
\end{proof}
\begin{remark}\label{double-dual-canonical}
For a finite abelian group, the evaluation map
$A\to\widehat{\widehat A}$, $a\mapsto(\chi\mapsto\chi(a))$, is
canonical, unlike the self-duality isomorphism above. Its bijectivity
is precisely the perfect evaluation pairing already proved in
\wikilink{s20--b006}{the finite duality theorem}, so that proof need
not be repeated. For a finite group all character values are roots
of unity. The isomorphism
$\mathbb Q/\mathbb Z\to\{\text{roots of unity in }\mathbb C^\times\}$,
$q+\mathbb Z\mapsto e^{2\pi iq}$, identifies the two dual notations
used here and in the source.
\end{remark}
